% !TeX program = xelatex
\documentclass[tikz,12pt]{standalone}
\usepackage[UTF8,scheme=plain,fontset=windows]{ctex}
\usepackage{tikz}
\usetikzlibrary{positioning,arrows.meta,calc}
\begin{document}
\begin{tikzpicture}[
  >=Latex,
  node distance=7mm and 8mm,
  every node/.style={font=\small},
  box/.style={draw,rounded corners,minimum width=20mm,minimum height=8mm,align=center},
  note/.style={align=center}
]
  \node[note] (title1) at (0,1.2) {状态侧};
  \node[note] (title2) at (0,-1.2) {测量侧};
  \node[box] (s1) at (2.0,1.2) {原子-腔发射};
  \node[box] (s2) at (5.0,1.2) {QFC};
  \node[box] (s3) at (8.2,1.2) {滤波记忆};
  \node[box] (s4) at (11.6,1.2) {输入端条件态};

  \draw[->] (s1)--(s2);
  \draw[->] (s2)--(s3);
  \draw[->] (s3)--(s4);

  \node[box] (m1) at (3.0,-1.2) {光纤损耗/相位};
  \node[box] (m2) at (6.6,-1.2) {分束器网络};
  \node[box] (m3) at (10.1,-1.2) {探测器 POVM};
  \node[box] (m4) at (13.5,-1.2) {等效 effect 推回};

  \draw[->] (m1)--(m2);
  \draw[->] (m2)--(m3);
  \draw[->] (m3)--(m4);

  \draw[->,dashed] (m4.north) -- (s4.south) node[midway,right] {\small $\Phi^{\dagger}$};
\end{tikzpicture}
\end{document}
